% --- ХУРААНГУЙ ХУУДАС ---
\newpage
\phantomsection
\addcontentsline{toc}{chapter}{Хураангуй}
\begin{center}
    \textbf{\large \MakeUppercase{Сувиллын захиалгын онлайн систем}}\\[0.2cm]
\end{center}

\vspace{0.5cm}

\begin{flushright}
    \begin{spacing}{1.2}
        {\normalsize Компьютерын ухааны тэнхим}\\
        {\normalsize Оюутан: Б.Зөнсод s21c118b}\\
        {\normalsize Удирдсан багш: Б.Батчулуун}
    \end{spacing}
\end{flushright}

\vspace{0.5cm}

\begin{center}
    \textbf{\MakeUppercase{Хураангуй}}
\end{center}

\vspace{0.3cm}

Өнөөдрийн байдлаар Өвөрхангай аймгийн Хужирт сумын сувиллууд үйлчлүүлэгч бүртгэх, өрөө захиалах үйл явцыг утсаар ярьж захиалдаг уламжлалт аргаар хэрэгжүүлж байна.
Захиалга баталгаажуулах, амрагчийн мэдээлэл бүртгэх, тайлан гаргах зэрэг ажлуудыг цаасан бүртгэл болон хүний гараар гүйцэтгэх нь мэдээллийн алдагдал, захиалгын давхардал, ажилтны ачаалал нэмэгдэх зэрэг үйл ажиллагааны хүндрэлийг үүсгэж байна.

Иймд энэхүү судалгааны ажлын хүрээнд “Ананд Хужирт” сувиллын захиалга бүртгэх, амрагчийн мэдээллийг удирдах, захиалгын төлөв хянах, тайлан боловсруулах боломж бүхий веб суурьтай захиалгын системийг хөгжүүлэхийг зорив. Тус систем нь сувиллын захиалгын үйл явцыг цахимжуулж, амрагч болон сувиллын ажилтнуудад мэдээллийг шуурхай солилцох, захиалгыг төвлөрсөн байдлаар удирдах боломжийг бүрдүүлснээр үйлчилгээний чанар болон үйл ажиллагааны үр ашгийг нэмэгдүүлэхэд чиглэнэ.

\vspace{1cm}

\noindent \textbf{Түлхүүр үгс:} сувиллын захиалгын систем, веб аппликейшн, захиалга удирдлага, үйлчлүүлэгчийн бүртгэл, мэдээллийн систем

% --- ENGLISH ABSTRACT ---
\newpage
\phantomsection
\begin{center}
    \textbf{\large \MakeUppercase{Web-based booking and management system for a health resort}}\\[0.2cm]
\end{center}

\vspace{0.5cm}

\begin{flushright}
    \begin{spacing}{1.2}
        {\normalsize Department of Computer Science}\\
        {\normalsize Student: B.Zunsod s21c118b}\\
        {\normalsize Supervisor: B.Batchuluun}
    \end{spacing}
\end{flushright}

\vspace{0.5cm}

\begin{center}
    \textbf{\MakeUppercase{Abstract}}
\end{center}

\vspace{0.3cm}

At present, health resorts located in Khujirt, Uvurkhangai Province, rely on traditional telephone-based methods for managing guest registration and room reservations.
Confirming reservations, recording guest information, and generating reports are carried out manually using paper-based records, which leads to operational difficulties such as data loss, duplicate bookings, and increased workload for staff.

Therefore, this research aimed to develop a web-based booking system for the ``Anand Khujirt'' health resort, providing capabilities for reservation registration, guest information management, booking status monitoring, and report generation. The proposed system digitalizes the reservation process by providing a centralized platform for managing booking information and facilitating efficient communication between guests and resort staff. As a result, it contributes to improved service quality, enhanced operational efficiency, and more effective management of resort activities.
\vspace{1cm}

\noindent \textbf{Keywords:} health resort booking system, web application, reservation management, guest registration, information system
